% generated from Bio 2/rotor_assembly.tex
\chapter{The sandpile group as a coordinate system on genome reconstructions}\label{chap:Bio2}
\chaptermark{Bio2}
\begin{chapterabstract}
\noindent
The reconstructions of a circular genome from its $k$-mer spectrum are the Eulerian
circuits of its de Bruijn multigraph $G$, and by the BEST theorem they correspond
bijectively to pairs (spanning arborescence, choice of exit order at every branch vertex).
The sandpile group $\K(G)$ acts simply transitively on the arborescences by rotor-routing
(Holroyd, Levine, M\'esz\'aros, Peres, Propp and Wilson), whose proof of BEST is itself a
rotor walk. We put the two facts to a use we have not found in the literature: $\K(G)$ is
a \emph{torsor of coordinates on the global part of the reconstruction space}, one chip
placed on a $k$-mer is an explicit rearrangement carrying one candidate genome to another
with the same spectrum, and the difference of two candidates is a group element. We
exhibit the moves, verify the torsor structure and the bijection on assembly graphs by
exhaustive enumeration, and show that the classical uniform sampler of Eulerian circuits
--- known since Altschul and Erickson (1985) and implemented in uShuffle --- is the torsor
in Smith-normal-form coordinates, which is what makes the distance between two samples
computable. The centrepiece is a real genome: for bacteriophage $\varphi$X174 the $12$-mer
spectrum leaves exactly one global choice, $\K=\Z/2$; the two reconstructions are
exhibited, they are $A\,X\,B\,Y\,A\,Z\,B\,U$ and $A\,X\,B\,U\,A\,Z\,B\,Y$ for two interleaved
repeated $12$-mers, and a single chip exchanges the $433$ and $181$ bp segments. This is
the $\Z/2$ of the interleaving theorem of the companion paper, realised in the first
genome ever sequenced. Seven checks, exact except the two sampling rows.
\end{chapterabstract}

\section{Introduction}\label{Bio2:sec:intro}

A companion paper \cite{Bio1} showed that the number of reconstructions of a circular
sequence from its $k$-mer spectrum factors, by the BEST theorem, into a local part
$\prod_v(d^+(v)-1)!$ determined by branch statistics and a global part $|\K(G_k)|$, the
order of the sandpile group, which branch statistics do not determine; and that the group
structure of $\K(G_k)$ records repeat length, multiplicity and interleaving. This note is
about the group itself rather than its order: what does the group \emph{do} to the set of
reconstructions?

\paragraph{What is known.} Three literatures bear on the question, and we claim nothing in
any of them.
\begin{itemize}
\item \emph{Counting and sampling.} The BEST theorem \cite{AEdB,SmithTutte} counts Eulerian
circuits of a directed Eulerian graph as (arborescences)$\times\prod(d-1)!$, and since
arborescences can be sampled uniformly in polynomial time \cite{CMN,PW,Wilson}, so can
Eulerian circuits; Creed and Cryan state this explicitly \cite{CC:Bio2}. (For \emph{undirected}
graphs counting is $\#$P-complete \cite{BW}; the directed case, which is ours, is easy.)
Bioinformatics has used exactly this since 1985 to shuffle sequences preserving $k$-let
counts \cite{AE,KMUW,uShuffle}, and to count reconstructions from $k$-mer spectra
\cite{ABCS,KSP}.
\item \emph{Rotor-routing.} Holroyd et al.\ \cite{HLMPPW} prove that the sandpile group of a
digraph with a sink acts freely and transitively on the spanning arborescences toward the
sink (their \S3), and in their \S4 give a bijective proof of BEST by a rotor walk: an
Eulerian tour is exactly one full turn of every rotor (Cor.~4.10). The rotor-router model
itself is the ``Eulerian walkers'' model of Priezzhev et al.\ \cite{PDDK}.
\item \emph{Sandpile groups of de Bruijn graphs.} Levine \cite{Levine} computed
$\K(\mathrm{DB}_n)$ for the complete binary de Bruijn graph; Bidkhori--Kishore and
Chan--Hollmann--Pasechnik extended this \cite{BK,CHP}. None concern genome graphs.
\end{itemize}

\paragraph{What is here.} We compose these into a statement about assembly
(Theorem~\ref{Bio2:thm:torsor}): the set of reconstructions, modulo the local exit choices, is
a torsor for $\K(G)$. Consequences: a chip placed on a $k$-mer is an explicit
rearrangement between two candidate genomes (Section~\ref{Bio2:sec:moves}); the difference of
two candidates is a group element, computable as a chip word; and the classical uniform
sampler is the torsor written in Smith-normal-form coordinates
(Section~\ref{Bio2:sec:sampling}). Section~\ref{Bio2:sec:phix} applies this to $\varphi$X174, where
the $12$-mer spectrum leaves one binary choice and we show what that choice is.

\paragraph{What we do not claim.} Not the torsor theorem, which is \cite[\S3]{HLMPPW}; not
the bijection between tours and (arborescence, rotor order), which is \cite[Cor.~4.10]{HLMPPW}
and which our check~(B) merely re-implements; not uniform sampling, which is
\cite{CMN,PW,CC:Bio2} and, in bioinformatics, \cite{AE,KMUW,uShuffle}. The contribution is the
reading of these as a coordinate system on assemblies, the identification of chip-firing
with rearrangement, and the $\varphi$X174 computation.

\section{Setting}\label{Bio2:sec:setting}

$S$ is a circular sequence, $G=G_k(S)$ its de Bruijn multigraph (vertices the distinct
$k$-mers, one arc per $(k{+}1)$-mer occurrence), which is connected and balanced with $S$
an Eulerian circuit; $\Rec(G)$ denotes the set of arc-rooted Eulerian circuits, i.e.\ the
reconstructions \cite[Rem.~1]{Bio1}. Fix a sink $s\in V$ and let $\Arb(G)$ be the set of
spanning arborescences oriented toward $s$; for balanced $G$ the count $|\Arb(G)|$ and the
group $\K(G)=\coker\Delta^{\mathsf T}$ do not depend on $s$ \cite[Lem.~4.12]{HLMPPW}.
Here $\Delta=D^{+}-A$ is the out-Laplacian with the row and column of $s$ deleted; its
determinant counts arborescences \emph{toward} $s$ because each non-sink vertex chooses
one of its out-arcs, and the transpose enters because firing $v$ sends one chip along
each out-arc, $\sigma\mapsto\sigma-\Delta^{\mathsf T}e_v$, so the firing lattice is the
column span of $\Delta^{\mathsf T}$. For balanced $G$, $\Delta^{\mathsf T}=D^{+}-A^{\mathsf T}$
is the in-Laplacian, so the choice is between two presentations of one group. A
\emph{rotor configuration} is a choice of one out-arc at every non-sink vertex; the
acyclic ones are exactly $\Arb(G)$. A chip placed at $v$ is routed to the sink by the
rule ``advance the rotor at the current vertex by one, then move along it''
\cite[\S3]{HLMPPW}; the map $\Arb(G)\to\Arb(G)$ this induces is the action of the class of
$e_v$ in $\K(G)$. The same action can be realised algebraically, by adding $e_v$ to a
recurrent chip configuration and toppling back to the recurrent representative, since
the rotor-router group and the sandpile group coincide \cite[Lem.~3.19]{HLMPPW}; we use
the dynamical description because it delivers the new arborescence directly.

By the contraction lemma of \cite[Lem.~4]{Bio1}, both $\K(G)$ and, arc for arc, $\Arb(G)$
are unchanged when every vertex of in- and out-degree $1$ is smoothed away, so all of the
following may be computed on the compacted (unitig) graph; the contracted arcs carry the
concatenated labels and a circuit on the compacted graph is spelled back into a
full-length sequence.

\section{The torsor}\label{Bio2:sec:torsor}

\begin{theorem}[Coordinates on reconstructions]\label{Bio2:thm:torsor}
Let $G$ be the de Bruijn multigraph of a circular sequence. Then
\begin{enumerate}
\item[(i)] $\K(G)$ acts freely and transitively on $\Arb(G)$ by rotor-routing
{\rm\cite[Lem.~3.12, 3.13, 3.17]{HLMPPW}};
\item[(ii)] the map $(T,\rho)\mapsto$ ``walk from the sink taking, at each vertex, the
next unused out-arc in the cyclic order $\rho_v$ with the tree arc of $T$ last'' is a
bijection from pairs (arborescence, exit orders) onto $\Rec(G)$
{\rm\cite[Cor.~4.10]{HLMPPW}};
\item[(iii)] consequently $\Rec(G)/\!\sim$, the reconstructions modulo the local exit
choices, is a $\K(G)$-torsor: after fixing one reconstruction, every other one is labelled
by a unique group element, and $|\Rec(G)|=|\K(G)|\cdot\prod_v(d^+(v)-1)!$.
\end{enumerate}
\end{theorem}

The theorem is a composition of two results of \cite{HLMPPW}; what we add is (iii) as a
statement about genomes, and the two verifications of Section~\ref{Bio2:sec:battery}: check~(A)
enumerates $\Arb(G)$ for six assembly graphs, builds the permutation representation of the
chip generators, and confirms that the generated group has order $|\K(G)|$, a single orbit,
and no fixed points; check~(B) implements the walk of (ii) and confirms it is a bijection
onto the brute-force enumeration of Eulerian circuits ($288$, $576$, $144$, $8$ circuits on
four graphs).

\begin{remark}[Why the quotient]\label{Bio2:rem:quotient}
The local choices at a branch vertex of out-degree $d$ --- which of the $d-1$ non-tree
exits to take first --- are what the local factor $\prod(d-1)!$ counts, and they are
independent of one another and of the arborescence. They are the ambiguity that branch
statistics \emph{do} see. The torsor lives on what remains, which \cite{Bio1} showed is
invisible to branch statistics; Theorem~\ref{Bio2:thm:torsor} says it nevertheless has the
structure of a finite abelian group, given explicitly by the Smith normal form of the
out-Laplacian.
\end{remark}

\section{A chip is a rearrangement}\label{Bio2:sec:moves}

Fix exit orders once and for all (say, in arc order). Then each arborescence determines
one reconstruction, and adding a chip at a $k$-mer $v$ moves the arborescence and hence
the reconstruction. Check~(M) does this on $S=\texttt{ACACCAAAACCA}$, $k=2$, where
$\K(G)=\Z/2\oplus\Z/6$:
\begin{center}\ttfamily\small
\begin{tabular}{ll}
base reconstruction & AAAACCAACCAC\\
$+$ one chip at \texttt{CA} & ACCAAAACCAAC \quad(6 positions changed)\\
$+$ one chip at \texttt{CC} & ACCAAAACCAAC \quad(6 positions changed)\\
$+$ one chip at \texttt{AA} & AAAACCAACCAC \quad(unchanged: $[e_{\texttt{AA}}]=0$ in $\K$)\\
\end{tabular}
\end{center}
Each row has the $2$-mer spectrum of $S$. The chip at \texttt{AA} is trivial in the group
and moves nothing, which is the correct behaviour: not every $k$-mer carries a nonzero
coordinate.

\begin{proposition}[Distance]\label{Bio2:prop:distance}
For two reconstructions with the same exit orders there is a unique $g\in\K(G)$ carrying
one arborescence to the other. Writing $g$ as a sum of chip classes $[e_v]$ with
non-negative coefficients --- always possible in a finite group --- gives a sequence of
physically realisable moves converting the first genome into the second. The minimal
number of chips $\delta(T,T')$ is a \emph{directed} distance on $\Rec(G)/\!\sim$: a
quasi-metric, since $\delta(T,T')$ and $\delta(T',T)$ need not agree. Admitting chips of
both signs, i.e.\ the symmetric generating set $\{\pm[e_v]\}$, gives the Cayley word
metric of $\K(G)$.
\end{proposition}

This is the sense in which the group is a coordinate system rather than merely a count.
Check~(W) computes the word for two independently sampled reconstructions of a
$27$-element torsor: three chips.

One thing a chip is \emph{not} is a single transposition. The rotor walk re-threads every
branch vertex it passes on its way to the sink, so one chip can change the tree arc at
several junctions at once. On the compacted $\varphi$X174 graph at $k=11$, which has ten
branch vertices, the nine generators change $6,5,3,2,4,2,3,1,0$ tree arcs respectively
(check~(J)); only one of them is a single re-threading. The directed distance of
Proposition~\ref{Bio2:prop:distance} is therefore not the transposition or reversal distance
of comparative genomics, and we see no immediate inequality between them.

\section{Sampling, and what the group adds}\label{Bio2:sec:sampling}

Uniform sampling of $\Rec(G)$ is classical: sample an arborescence uniformly by a
loop-erased random walk \cite{Wilson,PW} or by unranking \cite{CMN}, then choose exit
orders uniformly \cite{CC:Bio2}; this is what \cite{AE,KMUW} do and what uShuffle implements
\cite{uShuffle}. Theorem~\ref{Bio2:thm:torsor} gives the same distribution in a different form:
with $\K(G)\cong\bigoplus\Z/d_i$ from the Smith normal form $U\Delta^{\mathsf T}V=D$, draw
$r_i$ uniformly mod $d_i$, take $U^{-1}r$ as a chip configuration, and route it from a
fixed base arborescence. Check~(U) draws $2700$ samples this way on a torsor of order $27$
and finds all $27$ arborescences with $\chi^{2}=31.2$ on $26$ degrees of freedom
($p=0.22$); Wilson's sampler on the same graph gives $\chi^{2}=22.3$ ($p=0.67$,
check~(W)). The two are the same
distribution. What the coordinate form adds is not speed --- Wilson's algorithm is faster
--- but the ability to say \emph{where} a sample is. A shuffler such as uShuffle returns
independent draws and nothing else: two of its outputs stand in no computable relation,
and the sample space has no structure beyond its cardinality. In the torsor two samples
differ by a computable group element, a sample can be moved by a prescribed chip, the
whole space can be enumerated by stepping through Smith coordinates, and a sample can be
conditioned on a prescribed value of some coordinates. What it does \emph{not} give is a
linear way to impose local constraints, as the next remark records.

\begin{remark}[Local constraints are not cosets]\label{Bio2:rem:cosets}
A long read that spans a branch vertex $v$ fixes which out-arc of $v$ is the tree arc,
selecting $S_{e}=\{T\in\Arb(G):T(v)=e\}$. If several such constraints could be combined
by linear algebra in $\K(G)$, each $S_e$, transported to the group through the torsor,
would have to be a coset of a subgroup. It is not: on the graph of \texttt{ACACCAAAACCA}
at $k=2$ ($|\K|=12$) two of the seven non-empty constraint sets are not cosets, on
\texttt{CCAACAACCAAC} at $k=3$ three of nine, on \texttt{AABABBAABABB} two of eight;
only the single-repeat graph \texttt{AABBAABBAABB} ($\K=(\Z/3)^{2}$) has every
constraint a coset (check~(N)). The group structure is global; the local structure of
arborescences is not linear in it, and a ``congruence solver'' for long-read
disambiguation would have to be a search over the torsor, not a system of congruences.
\end{remark}

\section{$\varphi$X174: one global choice}\label{Bio2:sec:phix}

The companion paper computed the sandpile spectrum of bacteriophage $\varphi$X174
(NC\_001422.1, $5386$ bp) and found $\K(G_{12})=\Z/2$ and $\K(G_{11})=\Z/2\oplus\Z/66$
\cite[Tab.~1]{Bio1}. Theorem~\ref{Bio2:thm:torsor} says that the $12$-mer spectrum of the phage
leaves exactly \emph{two} global reconstructions. Check~(F) exhibits them.

\begin{proposition}[$\varphi$X174 at $k=12$]\label{Bio2:prop:phix}
The compacted $12$-mer graph of $\varphi$X174 has two branch vertices, the repeated
$12$-mers
\[
A=\texttt{CTTCTGCCGTTT},\qquad B=\texttt{CGTCAAGGACTG},
\]
at positions $457$ and $3205$ ($A$) and $264$ and $2760$ ($B$) of NC\_001422.1. Reading
the reference from the second $A$, the genome is
\begin{align*}
&A\,X\,B\,U\,A\,Z\,B\,Y,\\
&X=3217..263\ (2433\ \text{nt, wrapping}),\qquad U=276..456\ (181\ \text{nt}),\\
&Z=469..2759\ (2291\ \text{nt}),\qquad\qquad\quad\ \ Y=2772..3204\ (433\ \text{nt}),
\end{align*}
an interleaved pair. Its two reconstructions, both of length $5386$ and both with the
phage's $12$-mer spectrum, are the reference $A\,X\,B\,U\,A\,Z\,B\,Y$ and
$A\,X\,B\,Y\,A\,Z\,B\,U$, the segments following the two copies of $B$ exchanged; one chip
at a branch vertex carries one to the other.
\end{proposition}

The coordinates account for the whole genome: $2433+181+2291+433+4\cdot12=5386$, the
four unique segments and the four occurrences of the two $12$-mers; the exchange
permutes segments, so both reconstructions have the same length and, segment by segment,
the same content.

This is precisely the situation of the interleaving theorem \cite[Thm.~6]{Bio1} at $s=1$:
a single repeat of multiplicity two leaves no circular ambiguity (its two threadings are
rotations of one another), so a surviving $\Z/2$ at $k=L$ must come from two interleaved
repeats, and the one bit it records is which of the two $B$-exits follows which $A$. The
first genome ever sequenced has exactly one such pair of $12$-mers, and the group finds it.
At $k=11$ the $132$ classes of $\Z/2\oplus\Z/66$ are the corresponding statement one
level down; we have not analysed their structure.

The biology of the one bit is worth stating, because it is not what one might hope. By
the feature table of NC\_001422.1, the exchanged segment $U$ ($276..456$) spans the
junction of genes $C$ ($133..393$) and $D$ ($390..848$), and $Y$ ($2772..3204$) spans the
junction of the spike genes $G$ ($2395..2922$) and $H$ ($2931..3917$); the four repeat
occurrences lie in $C$, $D$, $G$ and $H$ respectively. Exchanging a $C/D$ junction for a
$G/H$ junction breaks the reading frames on both sides, so the alternative reconstruction
is not a viable isomer of the phage. It is the misassembly that a $12$-mer assembler is at
risk of producing, and the group locates it exactly: one bit, two named $12$-mers, two
named segments.

\section{Verification}\label{Bio2:sec:battery}

The script \texttt{rotor\_assembly.py} runs the eight checks of Table~\ref{Bio2:tab:battery} in
about five seconds, all passing; its output is archived as
\texttt{rotor\_assembly\_output.txt}, and the $\varphi$X174 sequence is distributed with it.

\begin{table}[ht]
\centering\small
\begin{tabular}{clc}
\toprule
key & check & mode\\
\midrule
(A) & Theorem~\ref{Bio2:thm:torsor}(i) on six assembly graphs, $|\K|\le27$: free and transitive & exact\\
(B) & Theorem~\ref{Bio2:thm:torsor}(ii): the walk is a bijection onto the enumerated circuits & exact\\
(M) & Section~\ref{Bio2:sec:moves}: chips are spectrum-preserving rearrangements & exact\\
(U) & Smith-coordinate sampler uniform, $\chi^{2}=31.2$ on $26$ df & sampled\\
(W) & Wilson's sampler agrees; the chip word between two samples & sampled\\
(F) & Proposition~\ref{Bio2:prop:phix} & exact\\
(J) & junctions re-threaded per chip on $\varphi$X174 at $k=11$: $0$ to $6$ & exact\\
(N) & Remark~\ref{Bio2:rem:cosets}: constraint sets are not cosets on $3$ of $4$ graphs & exact\\
\bottomrule
\end{tabular}
\caption{The verification battery (\texttt{rotor\_assembly.py}), $8/8$.}
\label{Bio2:tab:battery}
\end{table}

\section{Scope}\label{Bio2:sec:scope}

\emph{Quoted.} The torsor and the rotor proof of BEST \cite{HLMPPW}; uniform sampling
\cite{CMN,PW,Wilson,CC:Bio2} and its bioinformatic use \cite{AE,KMUW,uShuffle}; counting
\cite{ABCS,KSP}; the spectrum of $\varphi$X174 and the interleaving theorem \cite{Bio1}.

\emph{New here, as far as we can determine.} Theorem~\ref{Bio2:thm:torsor}(iii) as a statement
about assemblies; chip-firing as rearrangement and the word metric of
Proposition~\ref{Bio2:prop:distance}; the coordinate form of the sampler; and
Proposition~\ref{Bio2:prop:phix}.

\emph{Open.} (a) \emph{Double strands.} Real assembly graphs are bidirected
\cite{KM,MB}. The reverse-complement map sends an arc $u\to v$ to $\overline v\to\overline u$,
reversing direction, so the passage from the $k$-mer digraph to the bidirected graph is not
a covering of digraphs and the covering theory of critical groups \cite{RT} does not apply
as it stands; a sandpile theory for bidirected de Bruijn graphs would have to start from
the signed or bidirected Laplacian \cite{Zaslavsky}, and we know of none. (b) The $132$
classes of $\varphi$X174 at $k=11$, and more generally what the invariant factors of a
real genome's group mean position by position. (c) Whether the directed distance of
Proposition~\ref{Bio2:prop:distance} relates to any edit distance between assemblies; since a
single chip can re-thread many junctions (Section~\ref{Bio2:sec:moves}), it is not the
transposition distance, and any relation would have to pass through the structure of
rotor walks on the compacted graph.

\section*{Provenance of this chapter}

Unrefereed preprint. Theorem~\ref{Bio2:thm:torsor} is a composition of published results and is
verified, not reproved; Proposition~\ref{Bio2:prop:distance} is immediate from it;
Proposition~\ref{Bio2:prop:phix} is an exact computation reproducible in seconds. We searched
for prior use of the sandpile torsor on genome graphs and for any critical-group treatment
of bidirected de Bruijn graphs, and found neither; a search cannot establish a negative.
